\documentclass[10pt]{article}
\usepackage[paperwidth=42cm, paperheight=29.7cm, margin=2cm, top=1.8cm, bottom=1.5cm]{geometry}
\usepackage{fontspec}
\usepackage{microtype} % Kerning, ligatures, optical margin alignment
\usepackage[brazil]{babel}
\usepackage{multicol}
\usepackage{amsmath}
\usepackage{amssymb}
\usepackage{tikz}
\usetikzlibrary{calc, positioning, arrows.meta, shapes.geometric, decorations.pathreplacing}

% === TIPOGRAFIA (NEJM Editorial) ===
\setmainfont{Palatino}[Ligatures=TeX]
\setsansfont{Optima}
\newfontfamily\headerfont{Optima}[LetterSpace=3]
\newfontfamily\titlefont{Optima}[LetterSpace=8]

% === PALETA SÓBRIA ===
\definecolor{navy}{RGB}{26, 54, 93}        % Azul-marinho profundo
\definecolor{darktext}{RGB}{20, 20, 20}     % Quase preto
\definecolor{midgrey}{RGB}{120, 120, 120}   % Cinza médio
\definecolor{lightgrey}{RGB}{235, 235, 235} % Separadores
\definecolor{rulered}{RGB}{153, 27, 27}     % Acento vermelho sóbrio
\definecolor{ruleblue}{RGB}{30, 64, 175}    % Acento azul

\usepackage{xcolor}
\color{darktext}

\pagestyle{empty}
\setlength{\columnsep}{1.2cm}
\setlength{\columnseprule}{0.4pt}

% === MACROS TIPOGRÁFICAS ===
\newcommand{\sectiontitle}[1]{%
    \vspace{0.6cm}%
    {\headerfont\large\bfseries\textcolor{navy}{\addfontfeatures{LetterSpace=5}\MakeUppercase{#1}}}\\[2pt]%
    {\color{navy}\rule{3cm}{1.2pt}}\\[6pt]%
}

\newcommand{\titem}[1]{%
    \noindent\hangindent=1.5em\hangafter=1%
    {\textcolor{navy}{\small$\blacktriangleright$}}\hspace{0.5em}#1\par\vspace{3pt}%
}

\begin{document}
\begin{center}
\scalebox{0.82}{%
\begin{minipage}{\textwidth}

% ==================================================================
% CABEÇALHO ESTILO JOURNAL
% ==================================================================
\noindent
\begin{minipage}[b]{0.65\textwidth}
    {\titlefont\fontsize{38}{42}\selectfont\textcolor{navy}{\MakeUppercase{Oximetria de Pulso}}}\\[6pt]
    {\headerfont\large\textcolor{midgrey}{Monitorização Contínua e Não-Invasiva da Saturação Arterial de Oxigênio}}
\end{minipage}\hfill
\begin{minipage}[b]{0.30\textwidth}
    \begin{flushright}
        {\headerfont\small\textcolor{midgrey}{ONE PAGER $\cdot$ ICU CRITICAL CARE}}\\[3pt]
        {\small\textit{Nick Mark, MD} $\cdot$ \textit{Tradução Profissional}}\\[2pt]
        {\footnotesize\textcolor{midgrey}{onepagericu.com}}
    \end{flushright}
\end{minipage}

\vspace{4pt}
{\color{navy}\rule{\textwidth}{2pt}}
\vspace{2pt}
{\color{lightgrey}\rule{\textwidth}{0.4pt}}

\vspace{0.5cm}

% ==================================================================
% BLOCO INTRODUTÓRIO (antes das colunas)
% ==================================================================
\noindent
\begin{minipage}[t]{0.55\textwidth}
    {\headerfont\normalsize\textcolor{navy}{\addfontfeatures{LetterSpace=4}PRINCÍPIO FISIOLÓGICO}}\\[4pt]
    A oximetria de pulso constitui a medição contínua da saturação de oxigênio da hemoglobina por espectrofotometria.
    O dispositivo explora o fato de que a hemoglobina \textbf{oxigenada} (OxyHb) e a \textbf{desoxigenada} (DeoxyHb)
    absorvem luz vermelha (660\,nm) e infravermelha (940\,nm) de forma diferente.
    O componente \textit{pulsátil} (AC) do sinal isola a absorção arterial da absorção dos tecidos (DC).
\end{minipage}\hfill
\begin{minipage}[t]{0.40\textwidth}
    \centering
    % GRÁFICO DE ABSORÇÃO (compacto, dentro do header)
    \begin{tikzpicture}[scale=0.72, font=\sffamily\small]
        \draw[->, darktext!60] (0,0) -- (6.2,0) node[right, scale=0.7] {nm};
        \draw[->, darktext!60] (0,0) -- (0,4.2) node[above, scale=0.7] {Absorção};

        % Curvas
        \draw[ultra thick, ruleblue!80] (0, 3.5) .. controls (2, 2.8) and (3, 2) .. (5.5, 1)
            node[right, scale=0.7, ruleblue] {\textbf{OxyHb}};
        \draw[ultra thick, rulered!80] (0, 1) .. controls (1, 2) and (2, 3) .. (3, 3)
            .. controls (4, 3) and (5, 2.2) .. (5.5, 2.2)
            node[right, scale=0.7, rulered] {\textbf{DeoxyHb}};

        % Linhas de referência
        \draw[dotted, navy, thick] (0.8, 0) -- (0.8, 4) node[above, scale=0.6] {\textbf{660}};
        \draw[dotted, navy, thick] (4.8, 0) -- (4.8, 4) node[above, scale=0.6] {\textbf{940}};
    \end{tikzpicture}
\end{minipage}

\vspace{0.5cm}

% ==================================================================
% CONTEÚDO PRINCIPAL EM 3 COLUNAS
% ==================================================================
\begin{multicols}{3}

\sectiontitle{Sensor e Posicionamento}

O sensor emite luz que atravessa o leito capilar pulsátil. O fotodetector oposto mede a intensidade residual.

\titem{\textbf{Dedo} --- o mais preciso. Usa transmissão direta. O polegar pode ser marginalmente superior.}
\titem{\textbf{Orelha e Testa} --- usam reflexão. Podem ser mais confiáveis em vasoconstrição ou hipotermia.}
\titem{Nenhum local é universalmente superior. Tente vários para identificar o melhor sinal no paciente.}

\sectiontitle{Cor da Pele e Esmalte}

A melanina absorve luz de forma diferente, podendo causar \textbf{superestimação} da SpO\textsubscript{2} real em pacientes de pele negra --- especialmente em hipoxemia.

\titem{Pacientes negros têm \textbf{3$\times$ mais risco} de hipoxemia oculta (SpO\textsubscript{2} aparentemente normal com SaO\textsubscript{2} real baixa).}
\titem{Esmaltes escuros (azul, verde, preto) interferem significativamente na leitura.}

\sectiontitle{Estados de Baixo Fluxo}

Choque, vasoconstrição severa e suporte circulatório mecânico (ECMO, LVAD) reduzem ou eliminam o componente pulsátil do sinal.

\titem{Fluxo \textit{não-pulsátil} torna a oximetria convencional inútil.}
\titem{A amplitude da onda pletismográfica correlaciona-se com a perfusão periférica.}

\sectiontitle{Hemoglobinas Anormais}

\titem{\textbf{Carboxi-Hb (CO):} Leitura falsamente \textit{normal} (90--100\%). O sensor não distingue CO de O\textsubscript{2}. Suspeitar em incêndios.}
\titem{\textbf{Meta-Hb:} Leitura ``travada'' em $\sim$85\% independente do valor real.}
\titem{\textbf{Sulfa-Hb:} Leitura espúria baixa. Paciente pode não estar hipóxico.}
\titem{\textbf{HbA1c $>$ 7:} Leve superestimativa da SpO\textsubscript{2}.}

A \textbf{co-oximetria de pulso} (4+ comprimentos de onda) pode medir MetHb e CoHb diretamente.

\sectiontitle{Precisão em Baixa SpO\textsubscript{2}}

A curva de calibração foi derivada de voluntários saudáveis. Abaixo de SpO\textsubscript{2} $<$ 75\%, a correlação com a SaO\textsubscript{2} real deteriora-se significativamente.

\sectiontitle{Hiperoxemia}

Hiperóxia é prejudicial (especialmente pós-parada cardíaca), mas a oximetria não diferencia PaO\textsubscript{2} normal de supra-normal quando SpO\textsubscript{2} = 100\%.

\titem{Alvo razoável: SpO\textsubscript{2} $\geq$ 94\% (evitar tanto hipóxia quanto hiperóxia).}

\sectiontitle{Atraso na Leitura}

A medição periférica reflete o estado \textbf{5--15 segundos no passado}. O atraso é maior em hipotensão e menor com probes centrais (testa/orelha).

\titem{A SpO\textsubscript{2} exibida pode continuar caindo mesmo \textit{após} intubação correta.}

\sectiontitle{Corantes e Medicamentos}

Azul de metileno e indocianina causam \textbf{quedas falsas transitórias} da SpO\textsubscript{2}.

\vspace{0.3cm}
{\color{navy}\rule{\linewidth}{1.2pt}}
\vspace{0.3cm}

\sectiontitle{Análise da Onda Pletismográfica}

A onda possui componentes \textit{sistólicos} e \textit{diastólicos}. Sua morfologia fornece informações sobre o tônus vascular.

\vspace{0.2cm}
\begin{center}
\begin{tikzpicture}[scale=0.85, font=\sffamily\footnotesize]
    \fill[lightgrey!60] (0,0) .. controls (0.2,0) and (0.6,2.2) .. (0.8,2.2)
        .. controls (1.0,2.2) and (1.2,1) .. (1.5,1)
        .. controls (1.8,1.4) and (2.5,0) .. (3,0) -- cycle;
    \draw[thick, navy] (0,0) .. controls (0.2,0) and (0.6,2.2) .. (0.8,2.2)
        .. controls (1.0,2.2) and (1.2,1) .. (1.5,1)
        .. controls (1.8,1.4) and (2.5,0) .. (3,0);
    \node[above, navy, scale=0.8, font=\headerfont\bfseries] at (0.8, 2.2) {Sístole};
    \node[right, midgrey, scale=0.75, font=\headerfont] at (1.5, 1) {Nó Dicrótico};
    \draw[dashed, midgrey!60] (0.8,0) -- (0.8,2.2);
    \draw[dashed, midgrey!60] (1.5,0) -- (1.5,1);
\end{tikzpicture}
\end{center}

\vspace{0.1cm}

{\headerfont\small\textcolor{navy}{NORMAL}}\\[3pt]
\begin{tikzpicture}[scale=0.55]
    \draw[thick, navy] (0, 0) .. controls (0.2, 0) and (0.4, 1) .. (0.5, 1) .. controls (0.6, 1) and (0.7, 0.4) .. (0.8, 0.4) .. controls (0.9, 0.8) and (1.5, 0) .. (2, 0);
    \draw[thick, navy] (2, 0) .. controls (2.2, 0) and (2.4, 1) .. (2.5, 1) .. controls (2.6, 1) and (2.7, 0.4) .. (2.8, 0.4) .. controls (2.9, 0.8) and (3.5, 0) .. (4, 0);
\end{tikzpicture}

\vspace{0.15cm}

{\headerfont\small\textcolor{navy}{SINAL FRACO}} {\footnotesize\textcolor{midgrey}{--- $\downarrow$ perfusão}}\\[3pt]
\begin{tikzpicture}[scale=0.55]
    \draw[thick, midgrey] (0, 0) .. controls (0.2, 0) and (0.4, 0.3) .. (0.5, 0.3) .. controls (0.6, 0.3) and (0.7, 0.1) .. (0.8, 0.1) .. controls (0.9, 0.2) and (1.5, 0) .. (2, 0);
    \draw[thick, midgrey] (2, 0) .. controls (2.2, 0) and (2.4, 0.3) .. (2.5, 0.3) .. controls (2.6, 0.3) and (2.7, 0.1) .. (2.8, 0.1) .. controls (2.9, 0.2) and (3.5, 0) .. (4, 0);
\end{tikzpicture}

\vspace{0.15cm}

{\headerfont\small\textcolor{navy}{VASOCONSTRIÇÃO}} {\footnotesize\textcolor{midgrey}{--- $\uparrow$ dicrótico}}\\[3pt]
\begin{tikzpicture}[scale=0.55]
    \draw[thick, rulered!70] (0, 0) .. controls (0.2, 0) and (0.4, 1) .. (0.5, 1) .. controls (0.6, 1) and (0.7, 0.7) .. (0.8, 0.7) .. controls (0.9, 0.9) and (1.5, 0) .. (2, 0);
    \draw[thick, rulered!70] (2, 0) .. controls (2.2, 0) and (2.4, 1) .. (2.5, 1) .. controls (2.6, 1) and (2.7, 0.7) .. (2.8, 0.7) .. controls (2.9, 0.9) and (3.5, 0) .. (4, 0);
\end{tikzpicture}

\vspace{0.15cm}

{\headerfont\small\textcolor{navy}{VASODILATAÇÃO}} {\footnotesize\textcolor{midgrey}{--- $\uparrow$ sistólico}}\\[3pt]
\begin{tikzpicture}[scale=0.55]
    \draw[thick, ruleblue!70] (0, 0) .. controls (0.2, 0) and (0.4, 1.5) .. (0.5, 1.5) .. controls (0.6, 1.5) and (0.7, 0.2) .. (0.8, 0.2) .. controls (0.9, 0.5) and (1.5, 0) .. (2, 0);
    \draw[thick, ruleblue!70] (2, 0) .. controls (2.2, 0) and (2.4, 1.5) .. (2.5, 1.5) .. controls (2.6, 1.5) and (2.7, 0.2) .. (2.8, 0.2) .. controls (2.9, 0.5) and (3.5, 0) .. (4, 0);
\end{tikzpicture}

\vspace{0.4cm}

{\color{navy}\rule{\linewidth}{1.2pt}}
\vspace{0.2cm}

{\headerfont\small\textcolor{navy}{\addfontfeatures{LetterSpace=4}SUMÁRIO CLÍNICO}}\\[4pt]
{\small Sempre correlacione a onda pletismográfica com o pulso palpável.
Em dúvida (pele escura, intoxicação por CO, choque), a \textbf{gasometria arterial (SaO\textsubscript{2})} permanece o padrão-ouro.}

\end{multicols}

% ==================================================================
% RODAPÉ
% ==================================================================
\vfill
\noindent
{\color{lightgrey}\rule{\textwidth}{0.4pt}}\\[3pt]
\noindent
{\footnotesize\textcolor{midgrey}{\headerfont CRITICAL CARE TECHNICAL SERIES} \hfill \textit{Material exclusivamente educacional.} \hfill \headerfont\textcolor{midgrey}{PALS\,/\,ICU\,/\,2025}}

\end{minipage}%
}
\end{center}
\end{document}
